% GENERATED, edit the French unit or the translation table
\mxunit[7]{terminale-c10}{Antiderivatives and differential equations}\label{ch:7}

\mxsubtitle{S52, S53, S55, S58}

\begin{keypoint}[Objectives]Calculate antiderivatives; solve a differential equation with an initial condition; interpret an Euler approximation.\end{keypoint}

\section{Course notes}
\subsection{Antiderivatives on an interval}
At this stage we accept that every continuous function on an interval has antiderivatives; the link with area will be established in chapter 8.
An antiderivative of $f$ on an interval $I$ is a differentiable function $F$ such that $F'=f$ on $I$. If $F$ is an antiderivative, all the antiderivatives on $I$ are $F+C$, $C\in\R$.
\paragraph{Proof.} If $F'=G'=f$, then $(F-G)'=0$ on the interval; $F-G$ is therefore constant. Conversely, $(F+C)'=f$. The interval hypothesis is essential: on a disconnected domain, the constants may differ from one interval to another.
\begin{keypoint}[Table of antiderivatives --- to be read as the reverse of derivatives]
\[\begin{array}{c|c}
f(x)&F(x)\\\hline
x^n\ (n\in\Z,\ n\ne-1)&\frac{x^{n+1}}{n+1}\\
\frac{1}{x}&\ln|x|\quad\text{on an interval avoiding }0\\
\frac{1}{\sqrt{x}}&2\sqrt{x}\quad(x>0)\\
e^x&e^x\\
\cos x&\sin x\\
\sin x&-\cos x
\end{array}\]
The additive constant can always be added. The last two rows will be used again after the study of trigonometry.
\end{keypoint}
\subsection{Recognising a composition}
For a differentiable function $u$, the antiderivatives of $u'u^n$, $u'e^u$, $\frac{u'}{u}$ are respectively $\frac{u^{n+1}}{n+1}$ ($n\ne-1$), $e^u$, $\ln|u|$, on the intervals where the expressions make sense. Linearity allows us to handle sums.
\paragraph{Example.} An antiderivative of $6x(x^2+1)^2$ is $(x^2+1)^3$; an antiderivative of $e^{3x}$ is $\frac{e^{3x}}{3}$. To impose $F(0)=2$ in the latter case, take $F(x)=\frac{e^{3x}}{3}+\frac{5}{3}$.
\subsection{Differential equations}
A differential equation relates an unknown function to its derivatives. A solution must be checked on a specified interval.
\begin{keypoint}[Standard solutions]
$y'=ay$ has solutions $y(x)=Ce^{ax}$. For $a\ne0$, $y'=ay+b$ has solutions
\[y(x)=Ce^{ax}-\frac ba.\]
If $a=0$, the equation becomes $y'=b$ and $y(x)=bx+C$.
If $g$ is a known particular solution of $y'=ay+f(x)$, the solutions are $y=g+Ce^{ax}$.
\end{keypoint}
\paragraph{Proof for $y'=ay$.} If $y$ is a solution, $z(x)=e^{-ax}y(x)$ satisfies $z'=e^{-ax}(y'-ay)=0$; $z=C$. Conversely, the derivative of $Ce^{ax}$ is $aCe^{ax}$. For $y'=ay+b$, subtracting the constant solution $-\frac{b}{a}$ reduces to the homogeneous equation. For a variable right-hand side, do the same with $g$.
\paragraph{Example.} $y'=-2y+6$ and $y(0)=1$ give $y=3+Ce^{-2x}$, then $C=-2$. The solution $3-2e^{-2x}$ increases towards the equilibrium $3$ as $x\to+\infty$.
\subsection{Euler's method}
For $y'=ay+b$, with a step $h>0$, we approximate $y(t_k)$ by $Y_k$ according to $t_{k+1}=t_k+h$, $Y_{k+1}=Y_k+h(aY_k+b)$. This is an approximation using the tangent, not an equality on the exact solution. A smaller step can improve the approximation; check the error numerically on a known solution.


\mxnewpage
\section{Research activity}\textbf{Cooling of a drink --- 40 min.}\par A drink initially at $80\,^{\circ}\mathrm C$ is placed in a room kept at $20\,^{\circ}\mathrm C$. For $t\ge0$, let $T(t)$ be its temperature in degrees Celsius, where $t$ is the elapsed time in minutes. We accept the model $T' (t)=-\num{0,1}(T(t)-20)$ and $T(0)=80$.
\begin{enumerate}\item Solve this initial value problem and give the expression of $T(t)$.
\item Study the direction of variation of $T$ and find its limit at $+\infty$. Interpret the results.
\item Find the first instant at which the temperature is $40\,^{\circ}\mathrm C$. Give the exact value, then an approximate value in minutes.
\item Apply two steps of Euler's method with a step of $1$ minute. Compare the approximation obtained after two minutes with the exact value $T(2)$.
\item Give a concrete situation in which the assumption of a constant room temperature would no longer be appropriate.\end{enumerate}

\mxnewpage
\section{Progressive exercises}

\mxexerciseitem{1}{Practice}Find all the antiderivatives on $\R$ of the function $x\mapsto3x^2-4x+1$.

\mxexerciseitem{2}{Practice}Find the antiderivative $F$ of the function $x\mapsto\frac{2x}{x^2+4}$ on $\R$ that satisfies $F(0)=0$.

\mxexerciseitem{3}{Practice}Find the solution on $\R$ of the differential equation $y' =3y$ that satisfies $y(1)=2$.

\mxexerciseitem{4}{Practice}Find the solution on $\R$ of the differential equation $y' =-y+4$ that satisfies $y(0)=7$.

\mxexerciseitem{5}{Practice}Consider the differential equation $(E):y' =y-x$ on $\R$. Check that $g(x)=x+1$ is a solution of it, then find the solution of $(E)$ that satisfies $y(0)=0$.

\section{Textbook or past-paper work}\textbf{Barbazo, p. 228, exercises 25 and 30.} In S55: exercise 30, to bring out a constant solution. At home: exercise 25 in full, distinguishing a condition on $f$ from a condition on $f'$.

\mxnewpage
\section{Homework}\textbf{HW7 --- An evolution with variable input --- 60 min.}\par We look for a function $y$ defined and differentiable on $\intfo{0}{+\infty}$ such that $y' (t)=-y(t)+2t+1$ and $y(0)=3$.
\begin{enumerate}\item Find two real numbers $a$ and $b$ such that $g(t)=at+b$ is a solution of the differential equation.
\item Deduce the expression of the solution that satisfies $y(0)=3$.
\item Study the variations of this solution and state its minimum and the instant at which it is attained.
\item Write a Python function with argument $N$, a strictly positive integer, that applies Euler's method on $\intff{0}{2}$ with step $\frac2N$ and returns the approximation of $y(2)$. Compare with the exact value for $N=10$, then $N=100$.
\item Why is checking the differential equality at only a few instants not enough to prove that a function is a solution?\end{enumerate}

\ifmxteacher\mxnewpage\section{Teaching notes and solutions}\begin{teachernote}[Teacher notes]Distinguish the unknown function from the real variable. Have each antiderivative checked by differentiation. For $y'=ay+b$, look first for the constant solution. The case $a=0$ must remain visible. The trigonometric antiderivatives are given here as a reference table, but their assessed use awaits the trigonometry unit. For Euler, keep the units and compare with a known exact solution.\end{teachernote}

\subsection{Solutions to the activity}\begin{enumerate}
\item The equation is $T'=-\num{0,1}T+2$. A constant solution is $20$; all the solutions are $T(t)=20+Ce^{-\num{0,1}t}$. Since $T(0)=80$, $C=60$. We check $T'(t)=-6e^{-\num{0,1}t}=-\num{0,1}(T(t)-20)$ and $T(0)=80$.
\item $T'(t)=-6e^{-\num{0,1}t}<0$: the temperature decreases strictly. Since $e^{-\num{0,1}t}\to0$ at $+\infty$, $T(t)\to20\,^{\circ}\mathrm C$, while remaining above $20\,^{\circ}\mathrm C$.
\item $T(t)=40\Longleftrightarrow e^{-\num{0,1}t}=\frac{1}{3}$. Taking the logarithm, $-\num{0,1}t=-\ln3$, that is $t=10\ln3\simeq\num{10,9861}$ min. Strict decrease ensures that this is the only instant at which the temperature passes through $40\,^{\circ}\mathrm C$.
\item For $h=1$, $Y_{k+1}=Y_k-\num{0,1}(Y_k-20)$. Thus $Y_1=80-6=74$ and $Y_2=74-\num{5,4}=\num{68,6}$. The exact solution gives $T(2)=20+60e^{-\num{0,2}}\simeq\num{69,1238}$. Euler underestimates the temperature here by about $\num{0,5238}\,^{\circ}\mathrm C$.
\item The model assumes a room temperature of $20\,^{\circ}\mathrm C$ and a constant exchange coefficient of $\num{0,1}\,\mathrm{min}^{-1}$. A change of room, a lid or stirring can alter these assumptions; the formula is then no longer automatically appropriate.
\end{enumerate}

\paragraph{E1.}Linearity gives $F(x)=x^3-2x^2+x+C$, $C\in\R$. We check $F'(x)=3x^2-4x+1$. All the antiderivatives on the interval $\R$ differ by a constant.

\paragraph{E2.}The denominator is positive and its derivative is $2x$. The antiderivatives are $\ln(x^2+4)+C$. The condition at $0$ gives $C=-\ln4$, so $F(x)=\ln(x^2+4)-\ln4$; we check $F(0)=0$ and $F'(x)=\frac{2x}{x^2+4}$.

\paragraph{E3.}The solutions of $y'=3y$ are $Ce^{3x}$. The condition $Ce^3=2$ forces $C=2e^{-3}$. Thus $y(x)=2e^{3(x-1)}$, whose derivative is $3y(x)$ and whose value at $1$ is $2$.

\paragraph{E4.}The constant $4$ is a particular solution. The solutions are therefore $4+Ce^{-x}$. Since $4+C=7$, $C=3$. The solution is $y(x)=4+3e^{-x}$; its derivative $-3e^{-x}$ is indeed $-y+4$, and $y(0)=7$.

\paragraph{E5.}$g(x)=x+1$ satisfies $g'=1=g-x$. If $z=y-g$, then $z'=z$, so $y=x+1+Ce^x$. The condition $y(0)=0$ forces $C=-1$. We obtain $y=x+1-e^x$ and we check $y'=1-e^x=y-x$.

\subsection{Solutions to HW7}\begin{enumerate}
\item If $g(t)=at+b$, then $g'=a$ and $-g+2t+1=(2-a)t+1-b$. For the identity to hold for every $t$, we need $2-a=0$ and $a=1-b$. So $a=2$, $b=-1$, and $g(t)=2t-1$.
\item Let $z=y-g$. Subtracting the two equations, $z'=-z$, so $z=Ce^{-t}$. Thus $y(t)=2t-1+Ce^{-t}$. Since $y(0)=3$, $C=4$. The function $y(t)=2t-1+4e^{-t}$ does satisfy $y'=2-4e^{-t}=-y+2t+1$ and $y(0)=3$.
\item $y'(t)=2-4e^{-t}$. This derivative is negative for $t<\ln2$, zero at $\ln2$, positive afterwards. The solution decreases on $[0,\ln2]$ then increases on $[\ln2,+\infty[$. Its minimum is $y(\ln2)=2\ln2-1+2=1+2\ln2$.
\item 
\begin{lstlisting}
def euler_dm7(N):
    h = 2/N
    y = 3.0
    for k in range(N):
        t = k*h
        y = y+h*(-y+2*t+1)
    return y
\end{lstlisting}
The time used in the right-hand side is that at the start of the step. The program gives $Y_{10}\simeq\num{3,42950}$ and $Y_{100}\simeq\num{3,53048}$, whereas $y(2)=3+4e^{-2}\simeq\num{3,54134}$. The absolute errors are about $\num{0,11184}$ and $\num{0,01086}$.
To check these values without a loop, let $Z_k=Y_k-(2kh-1)$. Euler's recurrence gives $Z_{k+1}=(1-h)Z_k$ and $Z_0=4$. So $Y_N=3+4(1-\frac{2}{N})^N$.
\item A solution of a differential equation must satisfy an equality between functions on the whole interval, as well as the initial condition. A few numerical values only check a finite number of points; the symbolic substitution of question 2 is a check on the whole domain.
\end{enumerate}

\subsection{Solutions to the textbook or past-paper work}\paragraph{Exercise 30.} If $y$ is constant equal to $\ell$, the equation becomes $0=-\frac12\ell+3$, so $\ell=6$. Setting $z=y-6$, we obtain $z'=-\frac12z$; all the solutions are $y(x)=6+Ce^{-\frac x2}$, with $C\in\R$.
\paragraph{Exercise 25.}
\begin{enumerate}
\item $f'=\frac12f$ gives $f(x)=Ce^{\frac x2}$. Since $f(0)=1$, $C=1$.
\item $5f'=f$ gives $f(x)=Ce^{\frac x5}$. The condition $f(1)=0$ forces $Ce^{\frac15}=0$; since the exponential is never zero, $C=0$. The solution is the zero function.
\item $f'+f=0$ gives $f(x)=Ce^{-x}$, so $f'(x)=-Ce^{-x}$. The condition concerns the derivative: $f'(1)=2$ is equivalent to $-Ce^{-1}=2$, that is $C=-2e$. Thus $f(x)=-2e^{1-x}$.
\end{enumerate}

\fi

